% ============================================================
% EXAMEN DIAGNÓSTICO: ELECTROMAGNETISMO
% Curso de Oriente - CCH Oriente | Enero 2026
% Solo contenido — pegar en plantilla Overleaf
% ============================================================

\renewcommand{\tema}{Examen Diagnóstico — Electromagnetismo}

% ============================================================
\section*{EXAMEN DIAGNÓSTICO}
\subsection*{Electromagnetismo}
% ============================================================

% ============================================================
% PREGUNTAS
% ============================================================

\begin{preguntabox}
\subsection*{Pregunta 1}
Dos cuerpos cargados eléctricamente se atraen. ¿Qué podemos concluir sobre sus cargas?

\begin{enumerate}
    \item[A)] Ambas cargas son positivas.
    \item[B)] Ambas cargas son negativas.
    \item[C)] Las cargas son de signos opuestos.
    \item[D)] Las cargas son iguales en magnitud.
\end{enumerate}
\end{preguntabox}

\begin{preguntabox}
\subsection*{Pregunta 2}
¿En qué unidades se mide la resistencia eléctrica?

\begin{enumerate}
    \item[A)] Amperes (A)
    \item[B)] Volts (V)
    \item[C)] Watts (W)
    \item[D)] Ohms ($\Omega$)
\end{enumerate}
\end{preguntabox}

\begin{preguntabox}
\subsection*{Pregunta 3}
Una plancha eléctrica tiene una resistencia de 40 $\Omega$ y está conectada a una toma de 120 V. ¿Qué corriente circula por ella?

\begin{enumerate}
    \item[A)] 0.3 A
    \item[B)] 3 A
    \item[C)] 30 A
    \item[D)] 80 A
\end{enumerate}
\end{preguntabox}

\begin{preguntabox}
\subsection*{Pregunta 4}
Dos cargas eléctricas iguales de $3 \times 10^{-6}$ C cada una se encuentran separadas 0.3 m en el vacío. ¿Cuál es la magnitud de la fuerza entre ellas? ($k = 9 \times 10^9$ N$\cdot$m$^2$/C$^2$)

\begin{enumerate}
    \item[A)] 0.09 N
    \item[B)] 0.3 N
    \item[C)] 0.9 N
    \item[D)] 9 N
\end{enumerate}
\end{preguntabox}
\newpage
\begin{preguntabox}
\subsection*{Pregunta 5}
Una carga puntual de $2 \times 10^{-6}$ C genera un campo eléctrico. ¿Cuál es la magnitud del campo a 3 m de la carga? ($k = 9 \times 10^9$ N$\cdot$m$^2$/C$^2$)

\begin{enumerate}
    \item[A)] 500 N/C
    \item[B)] 1000 N/C
    \item[C)] 1800 N/C
    \item[D)] 2000 N/C
\end{enumerate}
\end{preguntabox}

\begin{preguntabox}
\subsection*{Pregunta 6}
Un calentador eléctrico tiene una resistencia de 30 $\Omega$ y opera con un voltaje de 60 V. ¿Cuál es su potencia?

\begin{enumerate}
    \item[A)] 2 W
    \item[B)] 90 W
    \item[C)] 120 W
    \item[D)] 1800 W
\end{enumerate}
\end{preguntabox}

\begin{preguntabox}
\subsection*{Pregunta 7}
Tres resistencias de 5 $\Omega$, 8 $\Omega$ y 7 $\Omega$ están conectadas en serie a una batería de 60 V. ¿Qué corriente circula por el circuito?

\begin{enumerate}
    \item[A)] 1 A
    \item[B)] 2 A
    \item[C)] 3 A
    \item[D)] 6 A
\end{enumerate}
\end{preguntabox}

\begin{preguntabox}
\subsection*{Pregunta 8}
Dos resistencias de 12 $\Omega$ y 6 $\Omega$ están conectadas en paralelo. ¿Cuál es la resistencia equivalente del conjunto?

\begin{enumerate}
    \item[A)] 2 $\Omega$
    \item[B)] 4 $\Omega$
    \item[C)] 9 $\Omega$
    \item[D)] 18 $\Omega$
\end{enumerate}
\end{preguntabox}

\begin{preguntabox}
\subsection*{Pregunta 9}
Dos condensadores de 6 $\mu$F y 3 $\mu$F se conectan en serie. ¿Cuál es la capacitancia equivalente del sistema?

\begin{enumerate}
    \item[A)] 1 $\mu$F
    \item[B)] 2 $\mu$F
    \item[C)] 4.5 $\mu$F
    \item[D)] 9 $\mu$F
\end{enumerate}
\end{preguntabox}

\begin{preguntabox}
\subsection*{Pregunta 10}
Una bobina de 100 vueltas experimenta un cambio de flujo magnético de 0.04 Wb en un intervalo de 0.02 s. ¿Cuál es la magnitud de la fem inducida?

\begin{enumerate}
    \item[A)] 4 V
    \item[B)] 20 V
    \item[C)] 200 V
    \item[D)] 2000 V
\end{enumerate}
\end{preguntabox}

% ============================================================
\newpage
\section*{RESPUESTAS Y RETROALIMENTACIÓN}
% ============================================================

\begin{respuestabox}
\subsection*{Pregunta 1}
\textcolor{verdecorrecto}{\textbf{Respuesta correcta: C) Las cargas son de signos opuestos.}}

\textbf{Justificación:}\\
La interacción entre cargas eléctricas sigue la regla fundamental: cargas de \textbf{signo opuesto se atraen}, cargas del \textbf{mismo signo se repelen}. Si los cuerpos se atraen, necesariamente una carga es positiva y la otra negativa.

\textbf{Respuestas incorrectas:}
\begin{itemize}
    \item \textbf{A) Ambas positivas:} Dos cargas positivas se \textit{repelerían}, no se atraerían.
    \item \textbf{B) Ambas negativas:} Dos cargas negativas también se \textit{repelerían}.
    \item \textbf{D) Iguales en magnitud:} La igualdad de magnitud no determina si hay atracción o repulsión; eso lo determina el \textit{signo}.
\end{itemize}
\end{respuestabox}

\begin{respuestabox}
\subsection*{Pregunta 2}
\textcolor{verdecorrecto}{\textbf{Respuesta correcta: D) Ohms ($\Omega$)}}

\textbf{Justificación:}\\
La \textbf{resistencia eléctrica} se mide en ohms ($\Omega$), en honor a Georg Simon Ohm. Las demás unidades corresponden a otras magnitudes eléctricas: el ampere mide corriente, el volt mide diferencia de potencial, y el watt mide potencia.

\textbf{Respuestas incorrectas:}
\begin{itemize}
    \item \textbf{A) Amperes (A):} Unidad de corriente eléctrica $I$.
    \item \textbf{B) Volts (V):} Unidad de diferencia de potencial (voltaje) $V$.
    \item \textbf{C) Watts (W):} Unidad de potencia eléctrica $P$.
\end{itemize}
\end{respuestabox}

\begin{respuestabox}
\subsection*{Pregunta 3}
\textcolor{verdecorrecto}{\textbf{Respuesta correcta: B) 3 A}}

\textbf{Justificación:}\\
\textbf{Fórmula utilizada:} Ley de Ohm despejada para la corriente:
$$I = \frac{V}{R}$$

\textbf{Sustitución:}
$$I = \frac{120 \text{ V}}{40 \text{ }\Omega} = 3 \text{ A}$$

\textbf{Respuestas incorrectas:}
\begin{itemize}
    \item \textbf{A) 0.3 A:} Error por invertir el cociente: $R/V = 40/120 = 0.33$ A — se dividió al revés.
    \item \textbf{C) 30 A:} Error de factor 10, posiblemente por confundir $R = 4\,\Omega$ con $R = 40\,\Omega$.
    \item \textbf{D) 80 A:} Resulta de sumar $V + R = 120 + 40$, lo cual no tiene fundamento físico.
\end{itemize}
\end{respuestabox}

\begin{respuestabox}
\subsection*{Pregunta 4}
\textcolor{verdecorrecto}{\textbf{Respuesta correcta: C) 0.9 N}}

\textbf{Justificación:}\\
\textbf{Fórmula utilizada:} Ley de Coulomb:
$$F = k \frac{|q_1| \cdot |q_2|}{r^2}$$

\textbf{Sustitución:}
$$F = (9 \times 10^9 \text{ N}\cdot\text{m}^2/\text{C}^2) \cdot \frac{(3 \times 10^{-6} \text{ C})(3 \times 10^{-6} \text{ C})}{(0.3 \text{ m})^2}$$
$$F = (9 \times 10^9) \cdot \frac{9 \times 10^{-12} \text{ C}^2}{9 \times 10^{-2} \text{ m}^2} = (9 \times 10^9) \cdot (10^{-10}) = 0.9 \text{ N}$$

\textbf{Respuestas incorrectas:}
\begin{itemize}
    \item \textbf{A) 0.09 N:} Error al calcular $r^2$: se usó $r^2 = 0.9$ en lugar de $r^2 = 0.09$ m$^2$.
    \item \textbf{B) 0.3 N:} Error por usar $r$ en lugar de $r^2$ en el denominador.
    \item \textbf{D) 9 N:} Se omitió dividir entre $r^2$, o se tomó $r = 0.1$ m en lugar de 0.3 m.
\end{itemize}
\end{respuestabox}

\begin{respuestabox}
\subsection*{Pregunta 5}
\textcolor{verdecorrecto}{\textbf{Respuesta correcta: D) 2000 N/C}}

\textbf{Justificación:}\\
\textbf{Fórmula utilizada:} Campo eléctrico de una carga puntual:
$$E = k \frac{|q|}{r^2}$$

\textbf{Sustitución:}
$$E = (9 \times 10^9 \text{ N}\cdot\text{m}^2/\text{C}^2) \cdot \frac{2 \times 10^{-6} \text{ C}}{(3 \text{ m})^2}$$
$$E = (9 \times 10^9) \cdot \frac{2 \times 10^{-6}}{9} = \frac{9 \times 10^9 \times 2 \times 10^{-6}}{9} = 2 \times 10^3 = 2000 \text{ N/C}$$

\textbf{Respuestas incorrectas:}
\begin{itemize}
    \item \textbf{A) 500 N/C:} Error por usar $r^2 = 9$ pero multiplicar por $k/4$ en lugar de $k$.
    \item \textbf{B) 1000 N/C:} Error por dividir el resultado correcto entre 2, quizás olvidando el factor $k$ completo.
    \item \textbf{C) 1800 N/C:} Error por usar $r = 3$ en lugar de $r^2 = 9$ en el denominador: $9\times10^9 \times 2\times10^{-6} / 3 = 6000/3 \neq 1800$.
\end{itemize}
\end{respuestabox}

\begin{respuestabox}
\subsection*{Pregunta 6}
\textcolor{verdecorrecto}{\textbf{Respuesta correcta: C) 120 W}}

\textbf{Justificación:}\\
\textbf{Fórmula utilizada:} Potencia eléctrica en función de voltaje y resistencia:
$$P = \frac{V^2}{R}$$

\textbf{Sustitución:}
$$P = \frac{(60 \text{ V})^2}{30 \text{ }\Omega} = \frac{3600 \text{ V}^2}{30 \text{ }\Omega} = 120 \text{ W}$$

\textbf{Respuestas incorrectas:}
\begin{itemize}
    \item \textbf{A) 2 W:} Error por dividir $V/R$ sin elevar al cuadrado y luego dividir de nuevo: confusión entre $P$ e $I$.
    \item \textbf{B) 90 W:} Error por calcular primero $I = 60/30 = 2$ A y luego usar $P = V \cdot R$ en lugar de $P = I^2 R$ o $P = VI$. Resulta $60 \times 30/20 \neq 90$ — es un error de fórmula no estándar.
    \item \textbf{D) 1800 W:} Error por calcular $P = V \times R = 60 \times 30$, confundiendo la fórmula con un producto directo sin el cuadrado ni la división.
\end{itemize}
\end{respuestabox}

\begin{respuestabox}
\subsection*{Pregunta 7}
\textcolor{verdecorrecto}{\textbf{Respuesta correcta: C) 3 A}}

\textbf{Justificación:}\\
\textbf{Paso 1 — Resistencia equivalente en serie:}
$$R_{eq} = R_1 + R_2 + R_3 = 5 \text{ }\Omega + 8 \text{ }\Omega + 7 \text{ }\Omega = 20 \text{ }\Omega$$

\textbf{Paso 2 — Corriente con Ley de Ohm:}
$$I = \frac{V}{R_{eq}} = \frac{60 \text{ V}}{20 \text{ }\Omega} = 3 \text{ A}$$

\textbf{Respuestas incorrectas:}
\begin{itemize}
    \item \textbf{A) 1 A:} Error por sumar solo dos de las tres resistencias ($5+8 = 13$ no da 1 A, o se confundió con otro dato).
    \item \textbf{B) 2 A:} Error por usar $R_{eq} = 30\,\Omega$, posiblemente al sumar $5+8+7+10$ con un valor imaginario.
    \item \textbf{D) 6 A:} Error por dividir el voltaje entre la resistencia más grande solamente: $60/7 \approx 8.6$, o por usar la mitad de $R_{eq}$: $60/10 = 6$.
\end{itemize}
\end{respuestabox}

\begin{respuestabox}
\subsection*{Pregunta 8}
\textcolor{verdecorrecto}{\textbf{Respuesta correcta: B) 4 $\Omega$}}

\textbf{Justificación:}\\
\textbf{Fórmula utilizada:} Resistencia equivalente de dos resistencias en paralelo:
$$R_{eq} = \frac{R_1 \cdot R_2}{R_1 + R_2}$$

\textbf{Sustitución:}
$$R_{eq} = \frac{12 \text{ }\Omega \times 6 \text{ }\Omega}{12 \text{ }\Omega + 6 \text{ }\Omega} = \frac{72 \text{ }\Omega^2}{18 \text{ }\Omega} = 4 \text{ }\Omega$$

\textbf{Respuestas incorrectas:}
\begin{itemize}
    \item \textbf{A) 2 $\Omega$:} Error por dividir entre el doble del denominador correcto: $72/36 = 2$.
    \item \textbf{C) 9 $\Omega$:} Error por promediar las resistencias: $(12+6)/2 = 9$ — el promedio no es la fórmula de paralelo.
    \item \textbf{D) 18 $\Omega$:} Error por sumar las resistencias como si fueran en \textit{serie}: $12+6 = 18$.
\end{itemize}
\end{respuestabox}

\begin{respuestabox}
\subsection*{Pregunta 9}
\textcolor{verdecorrecto}{\textbf{Respuesta correcta: B) 2 $\mu$F}}

\textbf{Justificación:}\\
\textbf{Fórmula utilizada:} Condensadores en serie:
$$\frac{1}{C_{eq}} = \frac{1}{C_1} + \frac{1}{C_2}$$

\textbf{Sustitución:}
$$\frac{1}{C_{eq}} = \frac{1}{6 \text{ }\mu\text{F}} + \frac{1}{3 \text{ }\mu\text{F}} = \frac{1}{6} + \frac{2}{6} = \frac{3}{6} = \frac{1}{2}$$
$$C_{eq} = 2 \text{ }\mu\text{F}$$

\textbf{Respuestas incorrectas:}
\begin{itemize}
    \item \textbf{A) 1 $\mu$F:} Error por calcular $1/C_{eq} = 1/6 + 1/6 = 1/3$, usando el mismo valor para ambos condensadores.
    \item \textbf{C) 4.5 $\mu$F:} Error por promediar: $(6+3)/2 = 4.5$ — el promedio no es la fórmula de condensadores en serie.
    \item \textbf{D) 9 $\mu$F:} Error por sumar directamente $C_1 + C_2 = 6 + 3$, que es la fórmula de condensadores en \textit{paralelo}, no en serie.
\end{itemize}
\end{respuestabox}

\begin{respuestabox}
\subsection*{Pregunta 10}
\textcolor{verdecorrecto}{\textbf{Respuesta correcta: C) 200 V}}

\textbf{Justificación:}\\
\textbf{Fórmula utilizada:} Ley de Faraday (magnitud):
$$|\mathcal{E}| = N \frac{|\Delta\Phi_B|}{\Delta t}$$

\textbf{Sustitución:}
$$|\mathcal{E}| = 100 \times \frac{0.04 \text{ Wb}}{0.02 \text{ s}} = 100 \times 2 \text{ V} = 200 \text{ V}$$

\textbf{Respuestas incorrectas:}
\begin{itemize}
    \item \textbf{A) 4 V:} Error por calcular solo $\Delta\Phi_B / \Delta t = 0.04/0.02 = 2$ V y luego dividir entre 2 en lugar de multiplicar por $N = 100$, o por omitir $N$ completamente y hacer $0.04/0.02/5$.
    \item \textbf{B) 20 V:} Error por multiplicar $N \times \Delta\Phi_B = 100 \times 0.04 = 4$ Wb y luego dividir entre $5\Delta t$ en lugar de $\Delta t$, o por usar $N = 10$ en lugar de $N = 100$.
    \item \textbf{D) 2000 V:} Error por confundir las potencias de diez: usar $\Delta t = 0.002$ s en lugar de 0.02 s, lo que multiplica el resultado por 10.
\end{itemize}
\end{respuestabox}
